\section{Pobrelma}
Oriunda do Mundo Antigo \cite{andreataAulaExpositivaPaulo2019}, a metodologia expositiva perseverou por incontáveis séculos, sobrevivendo a inúmeros paradigmas tecnológicos; perdurando até o presente contexto. É uma técnica na qual o emissor é protagonista no processo de ensino, enquanto o ouvinte - que devia receber maior atenção para a retenção do conteúdo - é relegado a coadjuvante nesse vital processo.

Essa problemática é constantemente levantada por acadêmicos, como é o caso de \cite{vasconcellosMetodologiaDialeticaEm1992} e \cite{pereiraCriticaMetodologiaTradicional2022}, que apontam o mero papel de ouvinte como prejudicial aos estudantes, que não são estimulados a engajar; tem uma postura meramente passiva. Eles ainda citam que o docente procura verificar se todos estão entendendo, mas por meio de perguntas vagas, indagando os estudantes quanto à sua compreensão. Nesse processo, os estudantes apenas concordam, e a conversa é finalizada no ato. Não há acompanhamento ativo, apenas uma estrutura hierarquizada e de efetividade discutível.

Os estudos de \cite{vasconcellosMetodologiaDialeticaEm1992} e \cite{pereiraCriticaMetodologiaTradicional2022} foram feitos em ambiente presencial de ensino, no  qual há mais facilidade de interação entre as partes, além de melhor controle sobre a atenção e engajamento da turma. Então, se mesmo nessas condições ambos os estudos se mostram céticos quanto à efetividade da aula expositiva, é de se esperar que em ambiente remoto esse cenário seja ainda pior.

Com a pandemia de COVID-19, o ensino à distância ganhou muito espaço, com destaque para o Ensino Superior. \cite{moranEnsinoSuperiorDistancia2009} identificou em seu trabalho um aumento exponencial nas matrículas por alunos em cursos à distância no Brasil; país relativamente pobre e com difícil acesso à tecnologia; mas que curiosamente tem boa parte de seus alunos universitários matriculados em cursos à distância, sejam eles 100\% ou semipresenciais.



% Muito Antiga
% Pouca Interação
% Centrada no Emissor
% Popularidade do Remoto
% Múltiplas Telas
% Dispersão de Foco